%----------------------------------------------------------------------------------------
%  "Not All Parameters Learn Alike"
%  Provisional acmart two-column stand-in for the JEDM acmtrans class.
%  Swap \documentclass to the JEDM class at submission; content is class-agnostic.
%  Compile: pdflatex -> bibtex -> pdflatex x2.
%----------------------------------------------------------------------------------------
\documentclass[sigconf,nonacm,screen]{acmart}

% --- strip ACM conference/copyright boilerplate ---------------------------------------
\settopmatter{printacmref=false}
\setcopyright{none}
\renewcommand\footnotetextcopyrightpermission[1]{}
\pagestyle{plain}
\settopmatter{printfolios=true}
\acmConference{}{}{}
\acmBooktitle{}
\setlength{\footskip}{12pt}

% --- packages -------------------------------------------------------------------------
\usepackage{booktabs}
\usepackage{amsmath}
\usepackage{graphicx}
\usepackage{tikz}
\usetikzlibrary{positioning,arrows,arrows.meta,backgrounds,calc}
\graphicspath{{figures/}}

\newcommand{\E}{\mathbb{E}}
\newcommand{\R}{\mathbb{R}}

\begin{document}

\title{Not All Parameters Learn Alike}
\subtitle{Parameter Recovery in Knowledge Tracing with Item Response Readouts}

\author{Wenrui Yuan}
\affiliation{%
  \institution{Section Cognition, Data and Education, University of Twente}
  \city{Enschede}
  \country{The Netherlands}
}
\email{}

\renewcommand{\shortauthors}{Yuan}

\begin{abstract}
Knowledge tracing models increasingly attach an item response readout so that a
prediction network also reports an ability, a discrimination, and a difficulty
for every learner and item. Practitioners then read those numbers. Trained only
to predict the next response, does the network actually recover them? We build a
modular tracer, a sequence encoder that swaps freely with an interpretable item
response decoder, and take parameter recovery against synthetic ground truth as
the yardstick. Across the encoder-by-decoder matrix, ability and difficulty come
back well, but the parameter that multiplies ability, the discrimination, comes
back worst. We trace the effect to a scale gauge. The logit multiplies
discrimination by ability, prediction loss constrains only that product, and
under a shared finite-width representation the split of scale between the slope
and the ability is under-determined, so the two compete. Decoupling the item
representation from the ability representation removes the competition, and the
additive difficulty, which never multiplies ability, never pays it. The effect is
a finite-data one that shrinks only in the joint limit of more training and more
learners. A second and separate penalty, a slow low-Fisher recovery rate, lands
on the generalized partial credit discrimination alone, and a state-conditioned
discrimination head is the targeted fix for it. The nominal response model
dissociates the two, since its slope multiplies ability and so trades off and is
rescued by decoupling, yet its Fisher information is not low, so the
state-conditioned head does nothing for it and even hurts. On EdNet and a KDD Cup
algebra set, discrimination reliability sits far below difficulty, static
decoupling can collapse it, and the state-conditioned head restores it. The
practice is to decouple always and to state-condition only the low-Fisher
parameter.
\end{abstract}

\keywords{knowledge tracing, item response theory, parameter recovery,
representation learning, learning dynamics, Fisher information, identifiability}

\maketitle

%======================================================================================
\section{Introduction}
%======================================================================================
Knowledge tracing follows a learner through a stream of graded responses and
predicts the next one \cite{corbett_knowledge_1994,piech_deep_2015}. A now common
design bolts an item response readout onto the prediction network, so that the
same model that forecasts the next answer also emits, at every step, an ability
$\theta_t$ for the learner and a discrimination $\alpha$ and a difficulty $\beta$
for each item \cite{yeung_deep-irt_2019,converse_irt_kt_2021}. The readout is
attractive because those three coordinates are readable. An instructor can look at
a discrimination and call an item sharp, look at a difficulty and call an item
hard, look at an ability trajectory and call a learner improving. The same three
numbers are increasingly load-bearing, since a difficulty and a discrimination are
what an item-selection rule consumes and a discrimination is the coordinate that
decides how much a response moves an ability estimate, so a wrong discrimination is
not a cosmetic blemish but a bad input to whatever the tracer feeds.

The reading rests on an assumption that is rarely tested. The network is trained
on one objective only, the likelihood of the observed responses. Nothing in that
objective supervises $\theta$, $\alpha$, or $\beta$ against a target. The item
response layer is a chosen parameterization of the output probabilities, and the
optimizer is free to reach the same predictions through any assignment of the
readout coordinates that produces them. Whether the recovered numbers track the
quantities they are named for is therefore an empirical question, and it is the
question this paper asks.

We ask it with a modular tracer. A sequence encoder, an LSTM, a memory network,
or a transformer, consumes the response history and produces a state. An item
response decoder, a binary two-parameter head, an ordinal generalized partial
credit head \cite{van_der_linden_generalized_1997}, or a nominal response head
\cite{bock_estimating_1972}, turns the state and an item embedding into a
distribution over the response categories. Encoder and decoder swap
independently, which lets us take one clean measurement, parameter recovery, and
read it across the whole grid. Recovery is the rank agreement between a recovered
parameter and its synthetic truth. It is the natural yardstick because the model
is identified only up to scale, so rank is what a scale-free estimator can hope to
pin \cite{baker_item_2004}.

The measurement returns an uneven verdict. Ability and difficulty recover well,
across encoders and across decoders. The discrimination recovers worst. The gap is
not an artifact of one backbone, because it survives every swap, and it is not
small, because it separates configurations that predict equally well.

The rest of the paper explains the verdict and acts on it. The parameter that
recovers worst is, in both the ordinal and the nominal decoder, the one that
multiplies ability inside the logit. That multiplication is a scale gauge. The
product of the slope and the ability is invariant when the slope is divided by a
constant and the ability multiplied by it, and prediction loss constrains only the
product. Under a shared representation of fixed width the division of scale between
the two is left to the optimizer, and the slope and the ability compete for it.
Giving the item its own representation, a step we call decoupling, ends the
competition, and the additive parameters, which never multiply ability, never
enter the competition at all. This coupling is a finite-data effect. It shrinks as
training and data grow together, and at any realistic budget it does not vanish.

A second penalty is layered on the ordinal discrimination and only on it. Its
Fisher information is low, because its leverage on the responses is concentrated
where responses are least informative, so it relaxes slowly and recovers late. The
fix for a low-Fisher parameter is to let it move with the state rather than sit as
a single item constant, a state-conditioned head. The nominal response model is the
control that keeps the two penalties apart. Its slope multiplies ability, so it
suffers the scale coupling and is helped by decoupling, yet its Fisher information
is not low, so the state-conditioned head does nothing for it and in fact hurts.
The ordinal discrimination has both problems, the nominal slope has only the first.

Our contribution is therefore twofold. First, a modular tracer and a recovery
benchmark across the encoder-by-decoder matrix, which establishes that the readouts
are unevenly recoverable and that the ordering is a property of the readout, not of
any one network. Second, a finite-data, finite-budget account of the ordering, in
which the parameter that multiplies ability recovers worst through a multiplicative
scale coupling that decoupling removes, kept distinct from a low-Fisher rate penalty
that a state-conditioned head removes, with the nominal response model as the
control that dissociates them. The two levers give one practice recommendation.
Decouple the item representation always, and state-condition only the parameter
whose Fisher information is low.

%======================================================================================
\section{Related work}
%======================================================================================
Knowledge tracing began with a latent mastery state updated by a hidden Markov
model \cite{corbett_knowledge_1994} and moved to recurrent networks with deep
knowledge tracing \cite{piech_deep_2015}, then to memory and attention backbones
\cite{zhang_dynamic_2017,pandey_self-attentive_2019,ghosh_context-aware_2020}. The
objective throughout is next-response prediction, and the item response readout is a
recent addition that reinterprets the prediction network's internals as measurement.
That reading began with Deep-IRT, which reads an ability and a difficulty out of a
memory network \cite{yeung_deep-irt_2019,zhang_dynamic_2017}, and continued through
explicit IRT-in-KT constructions \cite{converse_irt_kt_2021} and independent student
and item networks \cite{tsutsumi_state_2023,tsutsumi_deep_2024}. These works read the
recovered parameters as measurements. The closest neighbors to our finding saw the
symptom without the mechanism. \citet{converse_irt_kt_2021} report discrimination
recovering less well than the other parameters and read it as a data limitation,
and the beta4 item response model \cite{beta4irt_2023} reports an unexplained
shrinkage of discrimination magnitudes. We supply the mechanism those symptoms
share, a multiplicative scale coupling under a shared representation, and separate
it from a distinct low-Fisher rate penalty.

A parallel line recovers item response parameters by variational inference rather
than by prediction, from the variational item response model
\cite{wu_variational_2020} to variational temporal formulations \cite{vtirt_2023},
and \citet{ma_neural_2024} report the discrimination deficit under a variational
objective and attribute it to the variational gap. Our claim is narrower and about
a different training signal. Under plain prediction loss the deficit follows from
the scale gauge and from Fisher curvature, not from an inference gap, which is why
a variational method fit as well-specified generative inference can recover
discrimination best without contradicting us.

The nominal response readout has a precedent in option tracing, which predicts the
chosen distractor rather than only correctness \cite{ghosh_option_2021}. We use the
nominal response model \cite{bock_estimating_1972} not for its options but as a
control, because it gives each category its own slope and intercept and so splits a
slope on ability away from low Fisher information, a separation the generalized
partial credit model cannot make since it is the constrained nominal model with the
slopes tied to a single discrimination \cite{thissen_taxonomy_1986}. The rate half
of our account rests on standard learning-dynamics and information-geometry results,
the relaxation of modes near an optimum at rates set by curvature
\cite{saxe_exact_2014} and the Fisher metric of the natural gradient
\cite{amari_natural_1998}, and the gauge half restates the classical fact that item
response scales are fixed only up to an affine transform
\cite{baker_item_2004,bock_mle_1981}, read here as a property of a learned readout,
adjacent to the impossibility results for unsupervised disentanglement
\cite{locatello_disentanglement_2019,xi_indeterminacy_2023}.

%======================================================================================
\section{The modular tracer and the recovery benchmark}
%======================================================================================
\label{sec:tracer}

\input{fig_architecture.tex}

\subsection{The tracer}
The tracer has two swappable halves joined by a fixed interface (Figure
\ref{fig:architecture}). A sequence encoder maps the interaction history
$(q_1,r_1),\dots,(q_{t-1},r_{t-1})$ to a causal state $h_t$. From that state a
readout emits the person coordinate, the ability $\theta_t=\phi(h_t)$, while a
separate item embedding $e_i$ feeds the item coordinates. A decoder combines
$\theta_t$ and the item coordinates into a distribution over the $K$ response
categories, and the whole model is trained on the categorical likelihood of the
observed responses alone.

The encoder half is deliberately interchangeable. We use a plain LSTM, a key-value
memory network of the kind used in knowledge tracing \cite{zhang_dynamic_2017}, and
a transformer encoder \cite{vaswani_attention_2017,pandey_self-attentive_2019}. The
decoder half is likewise interchangeable across response formats. The binary
two-parameter head gives the logit $\alpha_i(\theta_t-\beta_i)$. The ordinal
generalized partial credit head \cite{van_der_linden_generalized_1997} accumulates
category logits $\sum_{m\le k}\alpha_i(\theta_t-\beta_{i,m})$, reducing to the
binary head at $K=2$. The nominal response head \cite{bock_estimating_1972} gives
each category its own slope and intercept, $a_{i,k}\theta_t+c_{i,k}$. All three
place a person term and an item term into one logit, which is the interface the
encoder and decoder agree on, and everything below concerns how that shared logit
distributes credit among its parts.

Two configuration choices are the levers the paper turns on, so we fix their meaning
here. The item representation is either \emph{shared} or \emph{decoupled}. A shared
item embedding of width $W$ feeds both the encoder input and the item readouts, tying
the ability's representation to the discrimination's, and the sweeps below vary $W$.
A decoupled item representation splits that in two, a narrow value embedding that
feeds the encoder and so shapes the ability, and a separate wide key embedding that
feeds the discrimination and difficulty readouts, so the two no longer share one
pool. The discrimination readout is either \emph{static} or \emph{state-conditioned}.
A static discrimination is one constant per item, read from the item key alone. A
state-conditioned discrimination is a function of the learner's current state as well,
$\alpha_{i,t}=\exp\bigl(w^\top[\,\text{key}_i;\,h_t\,]+b\bigr)$, so it varies along the
sequence rather than resting as a lookup. We train every configuration on the
categorical likelihood alone with Adam, the same schedule and the same $150$-epoch
stop, and read the readouts post hoc, so nothing but the response prediction shapes
what the parameters become.

\subsection{Recovery as the yardstick}
Because the likelihood fixes the item response scale only up to an affine transform
\cite{baker_item_2004}, no single-pass estimator can pin the raw magnitude of a
parameter, and a fair score reads rank rather than value. We take recovery to be
the sign-aligned Spearman rank correlation between a recovered parameter and its
generating truth, computed per item for the item parameters, averaged across
thresholds for the ordinal difficulties, and against per-learner ability for
$\theta$. Rank recovery near one means the model has ordered items and learners as
the truth did, which is what an interpretable readout must do to be read. Where
ground truth is absent, on real data, recovery is unavailable and we fall back to
split-half reliability, the rank agreement of a recovered parameter estimated on
two random halves of the learners, which needs no truth and so extends to the
field.

Rank is the primary score, but it hides where an estimator has settled on the raw
scale, so we read two diagnostics alongside it, the signed bias and the root mean
squared error of a recovered parameter against its truth. A rank near one with a
large bias says the ordering is right while the magnitude has drifted along the
scale gauge, which is the pattern the discrimination shows and the diagnostic that
first flagged it. Two reference points frame the recovery numbers. An offline
calibration of the same generalized partial credit model by marginal maximum
likelihood \cite{bock_mle_1981,chalmers_mirt_2012}, which sees the whole response
matrix at once rather than one causal pass, is the upper bound a single-pass tracer
is measured against, and a static per-learner ability table with no sequence model
is the non-sequential floor. The tracer sits between them, and the question is how
close to the offline bound each readout coordinate lands.

The synthetic generator is a generalized partial credit process with known
parameters, $N=800$ learners, $Q=60$ items, $K=4$ categories, sequences of length
$60$, ability $\theta\sim\mathcal{N}(0,1)$, discrimination
$\alpha\sim\mathrm{LogNormal}(0,0.5)$ as a moderate-spread prior consistent with
calibrated banks \cite{de_ayala_theory_2009}, and step thresholds drawn per item
around an item base. We train on prediction loss only, report at $150$ epochs, the
early-stopped operating point where a practitioner stops rather than a converged
one, and aggregate over eight seeds with $95\%$ confidence intervals unless noted.
This compact setting is the through-line for every result that follows. A per-item
recovery view on a larger instantiation ($N=5000$, $Q=200$) is shown for legibility
in Figure \ref{fig:recovery}.

\subsection{Recovery across the matrix}
\label{sec:matrix}
Table \ref{tab:matrix} reads recovery across the encoder-by-decoder grid. Two
facts hold in every cell. Ability and difficulty recover well, the ability rank
from $0.85$ to $0.97$ and the difficulty rank from $0.90$ to $0.99$, so the tracer
is genuinely modular and the readout is genuinely recoverable. The exception is
always the same coordinate. The discrimination, the coefficient that multiplies
ability, is the low and variable one, $0.71$ for the plain LSTM under an ordinal
head and $0.65$ for the binary head, and it climbs to about $0.89$ to $0.91$
exactly for the encoders that give the item its own wide key, the memory network,
the wide-key LSTM, and the transformer. The nominal slope $a_k$ recovers at $0.85$,
well above the ordinal discrimination in the matched narrow setting. The ordering
is a property of the readout that survives every swap of the network, and it
already points at the mechanism, since the discrimination recovers better precisely
where the item stops sharing the ability's representation.

The encoders that recover the discrimination best also carry more parameters, from
about seven thousand for the plain LSTM to thirteen thousand for the memory network
and thirty-two thousand for the transformer, so the matrix alone cannot separate a
representational cause from raw capacity. The controlled sweep of Section
\ref{sec:escape} draws the line. There, widening a single shared embedding adds
capacity yet lowers the ability even as it lifts the discrimination, and a decoupled
model with fewer parameters than the widest shared one recovers both at once, so the
gain is about how the representation is divided and not about how much of it there
is. The matrix shows the ordering is robust to the backbone, and the sweep shows the
cause is the representation split.

\begin{table}[t]
\caption{Recovery across the encoder-by-decoder matrix (synthetic, $N=800$,
$Q=60$, $K=4$, $150$ epochs, sign-aligned Spearman rank against truth). Ability and
difficulty recover well in every cell; the discrimination coefficient (bold) is the
low and variable one, and it rises for encoders that give the item its own wide
key. GPCM and 2PL rows are three seeds, the NRM row eight. The discrimination
column reports $\alpha$ for GPCM and 2PL and the slope $a_k$ for NRM; the location
column reports $\beta$ for GPCM and 2PL and the intercept $c_k$ for NRM. The last
column is the parameter count, so the discrimination recovery can be read against
capacity.}
\label{tab:matrix}
\small
\begin{tabular}{@{}llcccr@{}}
\toprule
Encoder & Decoder & Ability & \textbf{Discrim.} & Location & Params \\
\midrule
LSTM              & GPCM & 0.967 & \textbf{0.713} & 0.977 & 7.0k \\
Memory (DKVMN)    & GPCM & 0.905 & \textbf{0.894} & 0.911 & 12.9k \\
LSTM, wide key    & GPCM & 0.928 & \textbf{0.911} & 0.952 & 14.9k \\
Transformer       & GPCM & 0.851 & \textbf{0.894} & 0.901 & 31.7k \\
\midrule
LSTM              & 2PL ($K{=}2$) & 0.875 & \textbf{0.649} & 0.986 & 6.9k \\
LSTM              & NRM  & 0.953 & \textbf{0.846} & 0.887 & 7.0k \\
\bottomrule
\end{tabular}
\end{table}

\begin{figure}[t]
\centering
\includegraphics[width=0.98\columnwidth]{recovery_k4_item.pdf}
\caption{Per-item recovery on a synthetic generalized partial credit set
($N=5000$, $Q=200$, $K=4$). Rows are four instantiations of the tracer, a
static-ability baseline, a gated-recurrent encoder, a memory encoder with a shared
readout, and a memory encoder with a decoupled item pathway. The left column plots
recovered against true discrimination, the remaining columns recovered against the
true step difficulties, and the dashed line is the identity. The difficulty columns
hug the identity, while the discrimination column is visibly the most dispersed,
the same ordering Table \ref{tab:matrix} reports numerically.}
\label{fig:recovery}
\end{figure}

%======================================================================================
\section{Method and theory}
%======================================================================================
\label{sec:theory}
This section explains why the coefficient that multiplies ability recovers worst.
The account has two parts that must be kept apart, a multiplicative scale coupling
that is about the representation and a low-Fisher rate that is about curvature. We
first write the gradients honestly, then read the two gauges the logit carries,
then locate each penalty, then state what is finite and what is structural.

\subsection{The gradients, written honestly}
Take one response through the binary head, where the algebra is cleanest. With
logit $z=\alpha(\theta-\beta)$, predicted probability $p=\sigma(z)$, and observed
$y\in\{0,1\}$, the cross-entropy loss and the sigmoid link give the residual
$r:=\partial L/\partial z=p-y$. The chain rule through $z$, with
$\partial_\theta z=\alpha$, $\partial_\beta z=-\alpha$, and
$\partial_\alpha z=\theta-\beta$, gives
\begin{equation}
\partial_\theta L=r\,\alpha,\qquad
\partial_\beta L=-r\,\alpha,\qquad
\partial_\alpha L=r\,(\theta-\beta).
\label{eq:gpcm-grad}
\end{equation}
The clean scalar $r=p-y$ is the binary, canonical-link case. The ordinal head is
the same construction with a vector residual $r_k=p_k-y_k$, so the ability gradient
becomes the discrimination times an expected-score residual and the difficulty
gradient a cumulative residual, while the structure that matters is already visible
in \eqref{eq:gpcm-grad}. Ability and difficulty are both scaled by the
discrimination and carry opposite signs, and the discrimination itself is driven by
the person-item separation $\theta-\beta$.

The nominal head exposes the same skeleton with the slopes untied. With category
logit $z_k=a_k\theta+c_k$ and residual $r_k=p_k-y_k$,
\begin{equation}
\partial_{a_k}L=r_k\,\theta,\qquad
\partial_{c_k}L=r_k.
\label{eq:nrm-grad}
\end{equation}
The slope $a_k$ is a slope on ability, the nominal analogue of the ordinal
discrimination, and its gradient carries the factor $\theta$, while the intercept
$c_k$ carries none.

\subsection{The two gauges}
The logit carries two invariances, and each governs a different pair of
coordinates.

The \emph{scale gauge} is multiplicative and it is the one that matters here. The
logit multiplies the slope by the ability, $\alpha\theta$ in the ordinal head and
$a_k\theta$ in the nominal head, and that product is unchanged under
\begin{equation}
\alpha\to\alpha/s,\qquad \theta\to s\,\theta,\qquad s>0.
\label{eq:scale-gauge}
\end{equation}
Prediction loss sees only the product, so it pins the product and leaves the factor
$s$ free. When the slope and the ability are read from one shared representation of
fixed width, the optimizer must divide a finite pool of representational scale
between them, and because the data do not constrain the division, the two compete
for it. That competition is the slope-versus-ability trade-off. Reading the slope
from a separate item representation removes the shared pool and so removes the
competition, which is the decoupling fix. The additive intercept $c_k$ never
multiplies ability, so it carries no scale gauge with the ability and pays no such
trade-off, and this is the contrast the nominal head lets us draw.

The \emph{location gauge} is additive and structural. The logit reads the
difference $\theta-\beta$, which is unchanged under $\theta\to\theta+c$,
$\beta\to\beta+c$. This threshold identifiability holds as long as the ability
readout does not see the current item. If the ability readout $\phi(h_t,e_i)$ is
allowed to read item $i$, it can absorb a per-item location shift, one flat
direction is added to the loss for every item, and the difficulty becomes
unidentified at every budget. We keep the ability readout item-blind, $\phi(h_t)$,
by construction, which fixes the location gauge and lets us study the scale gauge in
isolation. The location failure does not vanish with data, but it is a design choice
we make once, not the subject of the experiments below.

\subsection{Where each penalty lands}
Near a gauge-fixed optimum each representational direction relaxes as
$\dot\delta=-H\delta$, with $\delta_c(t)=\delta_c(0)\,e^{-\lambda_c t}$ and the rate
$\lambda_c$ set by that direction's Fisher curvature
\cite{saxe_exact_2014,amari_natural_1998}. For a readout coordinate $\xi$ the
curvature is $I(\xi)=(\partial\eta/\partial\xi)^\top\Sigma\,(\partial\eta/\partial\xi)$
with $\Sigma$ the response covariance. Reading \eqref{eq:gpcm-grad}, the ability and
difficulty curvatures are $\alpha^2 w$ with $w=p(1-p)$, not suppressed, so those
directions are fast. The discrimination curvature is
\begin{equation}
I(\alpha)=\E\!\left[w\,(\theta-\beta)^2\right],
\label{eq:fisher-alpha}
\end{equation}
which is small exactly where response mass concentrates, at $\theta\approx\beta$, so
the discrimination is the slow direction and recovers late. This low-Fisher rate is
a property of the ordinal discrimination in particular, because its leverage
$(\theta-\beta)^2$ centers on the response mode, where a response carries the least
information about how sharply the item separates learners. The ratio of the fast
curvature to the slow one, $\kappa=I(\theta)/I(\alpha)$, is the conditioning number
of the readout, and a shared-step-size optimizer needs on the order of $\kappa$ more
steps to relax the discrimination than to relax the ability. A more discriminating
bank, with larger $\alpha$, raises $\kappa$ and so widens the gap, which is why the
effect is easiest to see where items separate learners well.

The nominal slope is the decisive check. Its leverage is $\theta$ itself, centered
at $\theta=0$ away from the mode, so its Fisher information is not low. Measured on
our synthetic set the ratio of the slope curvature to the intercept curvature is
about $0.90$, not the several-fold ability-to-discrimination gap the ordinal head
shows. The nominal slope multiplies ability, so it lives in the scale gauge of
\eqref{eq:scale-gauge} and shares the shared-representation trade-off, but its
Fisher information is ordinary, so it carries no low-Fisher rate penalty. That is
the dissociation the whole account turns on. The ordinal discrimination is both
multiplicative and low Fisher, the nominal slope is multiplicative but not low
Fisher, and the two levers should therefore act on them differently.

\subsection{Finite versus asymptotic}
Both penalties are finite. The low-Fisher rate is a slow relaxation mode plus a
finite-sample variance term, and the recovery gap between the slow coordinate and
the fast ones behaves like $e^{-T/\kappa}+O(\kappa/N)$ in training steps $T$ and
learners $N$, with $\kappa$ the ratio of the fast curvature to the slow one, so it
vanishes only in the joint limit of both growing. The scale coupling is finite in
the same spirit, since it is about how a finite-width shared code divides scale, and
it is removed structurally by decoupling. Only the location gauge is structural, and
we fix it by construction as above.

Two honesty notes bound the account. The ordering is a statement about the
eigenmodes of the Fisher matrix, not about raw parameters read one at a time,
because the discrimination and difficulty curvatures are coupled off-diagonally, and
the slow eigenvector is discrimination-dominated only where the discrimination is
above roughly one. Below that the order can invert, and we state the inversion
rather than hide it. And the equality of the Gauss-Newton curvature with the Fisher
information holds at a well-specified zero-residual optimum, so the phrase
``vanishes asymptotically'' is a within-model statement. On misspecified real data
the two curvatures differ and a residual gap need not close, which is why our
real-data claims below are about reliability and calibration, not about recovery
against a truth that does not exist.

%======================================================================================
\section{Experiments}
%======================================================================================
\label{sec:exp}
The experiments walk the account from Section \ref{sec:theory} through the compact
synthetic setting and then onto real data. Recovery follows the readout structure,
the shared representation forces the trade-off and decoupling escapes it, the effect
is finite-data, the state-conditioned head fixes the low-Fisher rate, the nominal
response model dissociates the two penalties, and the reliability picture holds on
EdNet and a KDD Cup algebra set.

\subsection{Recovery follows the readout structure}
The matrix of Table \ref{tab:matrix} and the scatter of Figure \ref{fig:recovery}
already make the first point. The coordinate that multiplies ability is the one that
recovers worst, in every encoder and in both graded decoders, while ability and
difficulty recover well throughout. The ordering also tracks training time (Table
\ref{tab:speeds}). In the narrow shared setting the discrimination rank climbs slowly,
from $0.10$ at $25$ epochs to $0.25$ at $50$ to $0.66$ at $150$, while the ability is
already at $0.94$ by $25$ epochs and the difficulty is high from the start. The two
axes come apart in a way that matters for model selection. Prediction quality plateaus
by about epoch $50$, well before the discrimination has finished recovering, so a
model chosen by its held-out loss is chosen while its discrimination is still moving,
and recovery is a distinct axis from prediction rather than a proxy for it. This is
the slow low-Fisher direction relaxing under \eqref{eq:fisher-alpha}, seen as a
trajectory.

\begin{table}[t]
\caption{Recovery against training time in the narrow shared setting (synthetic
GPCM, eight seeds, Spearman against truth). Ability recovers by the first
checkpoint, and difficulty is likewise high from the start, while the discrimination
climbs slowly and is still moving at epoch $150$, well past the epoch near $50$ where
prediction quality has already plateaued.}
\label{tab:speeds}
\small
\begin{tabular}{@{}lcc@{}}
\toprule
Epoch & Ability & \textbf{Discrimination} \\
\midrule
25  & 0.942 & \textbf{0.102} \\
50  & 0.961 & \textbf{0.251} \\
150 & 0.968 & \textbf{0.663} \\
\bottomrule
\end{tabular}
\end{table}

\subsection{The shared representation forces the trade-off, decoupling escapes it}
\label{sec:escape}
Table \ref{tab:escape} isolates the scale coupling with a single knob, the width of
the shared embedding. Widening it lifts the discrimination, from $0.66$ at width $8$
to about $0.92$ by width $48$, but the same widening lowers the ability, from $0.97$
to $0.90$, and the difficulty stays flat throughout. The discrimination and the
ability move in opposition, which is the trade-off the scale gauge predicts, one
pool of scale divided between two claimants. No single shared width serves both.

Decoupling escapes the trade-off rather than moving along it. A narrow value
embedding that feeds the encoder, paired with a separate wide key that feeds the
item readouts, recovers the discrimination at $0.92$ and the ability at $0.96$ at
once, above the shared frontier on the ability axis with non-overlapping intervals.
To match the decoupled discrimination a shared model must widen to about width $48$,
and there its ability has fallen to $0.90$, well below the decoupled $0.96$. The
escape is a genuine improvement in the ability coordinate at fixed discrimination,
not a repackaging of the same frontier.

\begin{table}[t]
\caption{The shared-width trade-off and the decoupling escape (synthetic GPCM, eight
seeds, $150$ epochs, Spearman against truth). Within the shared family, widening
lifts discrimination but lowers ability while difficulty stays flat. Decoupling
(narrow value width $8$, wide key width $64$) reaches high discrimination and high
ability together, above the shared frontier on ability.}
\label{tab:escape}
\small
\begin{tabular}{@{}lcc@{}}
\toprule
Configuration & Discrimination & Ability \\
\midrule
Shared, width 8   & 0.663 & \textbf{0.968} \\
Shared, width 16  & 0.843 & 0.946 \\
Shared, width 32  & 0.913 & 0.914 \\
Shared, width 48  & 0.924 & 0.903 \\
Shared, width 96  & 0.914 & 0.903 \\
\midrule
Decoupled (8, 64) & \textbf{0.921} & \textbf{0.959} \\
\bottomrule
\end{tabular}
\end{table}

\subsection{The effect is finite-data}
The trade-off is a finite-data one, and the decoupling advantage is what a finite
budget buys. Adding learners at fixed compute does not wash it out. Sweeping the
learner count from $100$ to $1600$ at $K=4$, the decoupling advantage on
discrimination stays essentially flat, $0.23$ at $N=100$ and $0.28$ at $N=1600$,
rather than shrinking toward zero. This matches the account of Section
\ref{sec:theory}, where the gap closes only in the joint limit of training and data
growing together and so persists at every realistic operating point. More data alone
does not make it go away, which is why the representational fix is worth making
rather than waiting out.

\subsection{The state-conditioned head fixes the low-Fisher rate}
Decoupling addresses the scale coupling, not the low-Fisher rate. The targeted fix
for a low-Fisher parameter is to let it move with the learner's state instead of
resting as one item constant, a state-conditioned discrimination head. A static
discrimination has to relax along its own slow, low-information direction, the one
\eqref{eq:fisher-alpha} shows is starved near the response mode. A
state-conditioned discrimination instead reads the learner's state, which is a
well-conditioned quantity that the ability has already driven to a good estimate
early in training, so the slope rides a fast-relaxing basis rather than waiting on
its own slow one. On the ordinal discrimination this helps, most clearly on real
data where the static decoupled discrimination can be starved, as Section
\ref{sec:real} shows, and it never lowers the ordinal discrimination in any cell we
ran, because a richer and faster source of variation for a slow coordinate can only
speed its relaxation. Whether that is a fix or a mere convenience is decided by the
nominal control, which we turn to now, since it is exactly where the
state-conditioned head stops helping.

\subsection{The nominal response model dissociates the two penalties}
\label{sec:nrm}
The nominal response model holds the multiplicative structure fixed while removing
the low Fisher information, and it splits the two levers cleanly (Table
\ref{tab:nrm}). The slope $a_k$ multiplies ability, so it shares the ability's
representation and pays the same trade-off, widening the shared embedding lifts
$a_k$ from $0.85$ to $0.94$ while dropping the ability from $0.80$ to $0.62$, the
signature of the scale gauge appearing even though $a_k$ is not low Fisher.
Decoupling clears it, the slope reaching $0.98$ and the intercept $0.98$ together,
above the best shared width at once. So far the nominal slope behaves like the
ordinal discrimination, which confirms that the trade-off is about sharing a
representation and not about low Fisher information.

The state-conditioned head is where they part. On the ordinal discrimination it
helps, but on the nominal slope it hurts, dropping $a_k$ from $0.98$ to $0.74$ while
the split-half reliability stays at $0.999$. The recovered slope is stable and
wrong, agreeing with itself across halves while disagreeing with the truth, because
the extra state variation is noise for a coordinate whose Fisher information is
already adequate. The state-conditioned head helps a parameter only where its Fisher
information is low, and the nominal slope, multiplicative but not low Fisher, is
helped by decoupling and hurt by state conditioning. The ordinal discrimination,
which is both, needs both levers, and the nominal slope shows that the levers are
genuinely two and not one. A note of caution belongs on the record. Decoupling the
nominal slope alone leaves its intercept starved at $0.39$, which is the nominal
head's own internal competition between its slope and its intercept for the item
key, a separate matter from the slope-versus-ability coupling, and we control it by
decoupling both.

\begin{table}[t]
\caption{The nominal response model dissociates the two penalties (synthetic, eight
seeds, $150$ epochs). Top, the slope $a_k$ multiplies ability and pays the
shared-representation trade-off just as the ordinal discrimination does, and
decoupling clears it. Bottom, state conditioning helps the low-Fisher ordinal
discrimination but hurts the nominal slope, whose Fisher information is not low (the
recovered slope is stable across halves yet far from truth).}
\label{tab:nrm}
\small
\begin{tabular}{@{}lccc@{}}
\toprule
Configuration & Slope $a_k$ & Intercept $c_k$ & Ability \\
\midrule
Shared, width 8   & 0.846 & 0.887 & 0.797 \\
Shared, width 32  & 0.939 & 0.895 & 0.620 \\
Decoupled         & \textbf{0.980} & \textbf{0.982} & 0.853 \\
\midrule
\multicolumn{4}{@{}l}{\emph{State-conditioned slope}} \\
Decoupled, static  & \textbf{0.980} & 0.982 & 0.853 \\
Decoupled, dynamic & \textbf{0.735} & 0.969 & 0.857 \\
\bottomrule
\end{tabular}
\end{table}

\subsection{Real data}
\label{sec:real}
On real data there is no generating truth, so we read split-half reliability rather
than recovery. Table \ref{tab:real} reports it on EdNet \cite{choi_ednet_2020} and a
KDD Cup algebra set \cite{stamper_kdd_2010}, both scored into three ordered
categories from correctness and effort. EdNet contributes $3000$ learners over $250$
items with a median sequence of $115$, the three categories being wrong, correct but
slow, and correct and fast. The algebra set contributes $2476$ learners over $250$
items with a median sequence of $172$, the three categories reading a correct first
attempt, a correct answer after no hint, and a correct answer that needed a hint.
Both pass a monotonicity check, the learner accuracy correlating with the mean
category at rank about $0.85$, so the ordinal coding is faithful. The synthetic
ordering survives.
Discrimination reliability sits far below difficulty on both sets, $0.77$ against
$0.94$ on EdNet and $0.62$ against $0.81$ on the algebra set, so the coordinate that
was hardest to recover against truth is also the least stable across learners in the
field.

The two levers behave as the synthetic account predicts. Static decoupling can
starve the discrimination, most sharply on the algebra set where it collapses to
$0.39$, because decoupling removes the shared pool but leaves a low-Fisher
coordinate to relax on its own. The state-conditioned head rescues it, lifting the
algebra discrimination from $0.39$ to $0.64$ and the EdNet discrimination from
$0.64$ to $0.78$, and it never lowers discrimination reliability in any cell. The
practice that emerges is to decouple and to state-condition the discrimination
together, since decoupling clears the scale coupling and the state-conditioned head
covers the low-Fisher rate that decoupling leaves exposed.

The nominal readout on real data adds a coverage caveat that the reliability lens
makes visible. On EdNet option tracing, where each of more than ten thousand items
is seen only about a dozen times, the decoupling escape reverses, shared ability
reliability at $0.68$ over decoupled at $0.29$, because a wide item key cannot be
estimated from a dozen observations. Aggregating to knowledge components, about
$845$ observations per unit, the escape revives, decoupled at $0.66$ against shared
at $0.70$, while the state-conditioned head still does not help the nominal slope,
$0.49$ below the static $0.66$. The portable finding is the one that does not depend
on coverage, the state-conditioned head helps the low-Fisher ordinal discrimination
and not the ordinary-Fisher nominal slope, exactly as the synthetic dissociation
says.

\begin{table}[t]
\caption{Real-data split-half reliability (rank agreement across two random halves
of the learners, four splits). Discrimination reliability sits far below difficulty
on both datasets. Static decoupling can starve the discrimination (bold on the
algebra set), and the state-conditioned head restores it without ever lowering it.}
\label{tab:real}
\small
\begin{tabular}{@{}llcc@{}}
\toprule
Dataset & Configuration & Discrimination & Difficulty \\
\midrule
EdNet   & Shared, static     & 0.771 & 0.940 \\
        & Shared, dynamic    & 0.790 & 0.941 \\
        & Decoupled, static  & 0.639 & 0.953 \\
        & Decoupled, dynamic & 0.784 & 0.955 \\
\midrule
Algebra & Shared, static     & 0.621 & 0.810 \\
        & Shared, dynamic    & 0.628 & 0.836 \\
        & Decoupled, static  & \textbf{0.391} & 0.888 \\
        & Decoupled, dynamic & 0.640 & 0.870 \\
\bottomrule
\end{tabular}
\end{table}

Real item-parameter geometry corroborates the reliability picture (Figure
\ref{fig:assist}). On the proxy-ordinal ASSISTments releases, which score effort
into ordered categories and carry no ground truth, the recovered discrimination
distribution across items is dispersed and heavy-tailed relative to the difficulty
distribution, the visual signature of the least stable coordinate, consistent with a
channel whose scale is only weakly pinned by the data. The 2009 release spans $2596$
learners over $3910$ items and the 2017 release $1699$ learners over $772$ items, and
in both the memory encoder with a decoupled item pathway spreads its discrimination
into a longer tail than the shared-readout encoders, which is what a coordinate free
to absorb loose scale does when nothing anchors its magnitude. The point is not that
one instantiation is right, since no truth adjudicates it, but that the coordinate
the synthetic and reliability results single out is exactly the coordinate whose
real-data geometry is least settled.

\begin{figure*}[t]
\centering
\begin{minipage}{0.48\textwidth}
\centering
\includegraphics[width=\linewidth]{assist2009_ord_item_params_3model.pdf}
\\{\small (a) ASSISTments 2009 ($Q=3{,}910$)}
\end{minipage}\hfill
\begin{minipage}{0.48\textwidth}
\centering
\includegraphics[width=\linewidth]{assistments_item_params_3model.pdf}
\\{\small (b) ASSISTments 2017 ($Q=772$)}
\end{minipage}
\caption{Per-item parameter geometry on the proxy-ordinal ASSISTments releases,
which carry no ground truth. Each panel shows, across items, the density of the
recovered discrimination (top) and of the recovered step difficulties (bottom) for
three instantiations of the tracer, a gated-recurrent encoder, a memory encoder with
a shared readout, and a memory encoder with a decoupled item pathway. The
discrimination density is dispersed and heavy-tailed relative to the difficulty
density, and the decoupled item pathway spreads it furthest, the geometry of the
least stable coordinate.}
\label{fig:assist}
\end{figure*}

%======================================================================================
\section{Discussion}
%======================================================================================
\subsection{Two levers, one recommendation}
The recovery deficit on the coordinate that multiplies ability has two causes, and
each has its own lever. The scale coupling under a shared representation is a
representational cause, and its lever is decoupling, a narrow value embedding for
the encoder and a separate wide key for the item readouts. The low Fisher
information of the ordinal discrimination is a curvature cause, and its lever is the
state-conditioned head. The two are not the same knob, which the nominal response
model proves by carrying the first without the second. The recommendation is
therefore specific. Decouple the item representation always, because any slope on
ability pays the shared-representation trade-off and a wide key clears it at no cost
to ability whether or not the slope's Fisher information is low. State-condition
only the parameter whose Fisher information is low, because the state-conditioned
head rescues the ordinal discrimination and degrades the nominal slope, so it is a
targeted fix and not a default.

\subsection{What a trustworthy discrimination buys}
The discrimination is the coordinate a reader is most tempted to act on, since it
labels an item sharp or dull, and it is the coordinate this paper finds least
trustworthy under prediction training. The levers matter because they make the label
safer to read. A decoupled, state-conditioned discrimination is the one
configuration that recovers the parameter against synthetic truth and holds its rank
across learners on real data, which is the minimum an interpretable readout owes the
instructor who reads it. The same estimate is what a downstream selection rule would
consume, so a discrimination one can trust is not only a cleaner number but a better
input to whatever the tracer feeds.

\subsection{Boundaries}
The account is sized to its evidence and we mark its edges. The synthetic results
use eight seeds at a single number of categories, and the ordering law is per
eigenmode, so below a discrimination of about one it can invert. The real-data leg
is reliability and calibration rather than recovery against a truth that does not
exist, and the nominal escape is coverage-contingent, reversing where items are seen
only a handful of times. Two confirmations remain to be run and we state each
plainly. A controlled ablation that holds Fisher information fixed while varying only
whether a parameter enters multiplicatively or additively with ability would turn
multiplicative entry from the best explanation into a demonstrated cause, and the
nominal slope and intercept already give a natural version of that contrast. An
oracle clamp that fixes the ability at its true value while the discrimination alone
relaxes would isolate the rate claim from the co-learning of the ability, and a knob
that varies the ability's Fisher information at fixed category count would pin the
rate coefficient rather than only its form and sign.

%======================================================================================
\section{Conclusion}
%======================================================================================
Prediction-trained knowledge tracing with an item response readout does recover
interpretable parameters, but not alike. Ability and difficulty come back well
across encoders and decoders, and the coefficient that multiplies ability comes back
worst. Two mechanisms govern the gap, and keeping them apart is the paper's point. A
multiplicative scale coupling under a shared representation makes any slope on
ability compete with the ability for a finite pool of scale, and decoupling the item
representation ends the competition. A low Fisher information makes the ordinal
discrimination relax slowly, and a state-conditioned head speeds it. The nominal
response model dissociates the two, since its slope is multiplicative but not low
Fisher, so decoupling helps it and state conditioning hurts it. The practice follows
without ambiguity. Decouple always, and state-condition only the parameter whose
Fisher information is low.

\bibliographystyle{ACM-Reference-Format}
\bibliography{ref}

\end{document}
